Before any of what follows, a profile is needed, because it changes which
number matters. On a 26-object scenario over four orbits the initial screen
costs \SI{14.7}{\second}, the SGP4 re-screen \SI{11.2}{\second}, latent-row
enumeration \SI{1.0}{\second}, and the linear program \SI{0.24}{\second}. The
LP is \textbf{1\%} of the runtime. Accelerating it, which is what the
predict-and-certify machinery does, is therefore worth at most one per cent
however well it works --- and the three changes that actually mattered were
exact algorithmic ones: memoising a rotation matrix per (object, epoch) and
writing out the cross products by hand (\textbf{2.1$\times$} on screening,
conjunction lists bit-identical), restricting the re-screen to pairs touching a
burn since nothing else moved (\textbf{2.5--14$\times$}, induced counts
identical), and the lazy closure below (up to \textbf{5$\times$} on the LP,
$\dv$ identical to 0.0e+00). End to end, \textbf{3.8$\times$}.

A fourth exact change removes the remaining scaling limit, which is memory
rather than time. The slack block is diagonal --- one slack column per resolve
or latent row --- so each latent row carries $n_{\text{lifted}}+1$ nonzeros
while the column count itself grows with the latent count, making dense storage
quadratic in the only dimension that grows. At \num{20840} rows the matrix is
\SI{0.23}{\percent} dense. Assembling coordinate triplets and materialising CSR
above four million elements takes peak assembly memory from \SI{3575}{\mega\byte}
to \SI{84}{\mega\byte}, a factor of \num{42.5}, with the matrix entrywise
identical, the solution vector bit-identical, and 56 of 56 benchmark rows
reproduced to \num{0.0e+00}.

With that said, the lazy-constraint closure returns the exact optimum for any
guess, and its speedup grows with the problem. With an oracle guess on the
\texttt{induced-cascade} family, varying the neighbour count and the latent
grid step:

\begin{center}
\begin{tabular}{rrrrr}
\toprule
latent rows & rows used & full solve & lazy solve & speedup \\
\midrule
   68 & 1 & \SI{4.5}{\milli\second}   & \SI{2.5}{\milli\second}  & $1.8\times$ \\
  358 & 2 & \SI{18.5}{\milli\second}  & \SI{5.4}{\milli\second}  & $3.5\times$ \\
  820 & 2 & \SI{48.5}{\milli\second}  & \SI{8.3}{\milli\second}  & $5.9\times$ \\
 1716 & 2 & \SI{145.8}{\milli\second} & \SI{16.1}{\milli\second} & $\mathbf{9.0\times}$ \\
\bottomrule
\end{tabular}
\end{center}

Objective gap exactly zero at every size. Exactness was verified against an
empty, a learned, and a deliberately adversarial guess over 27 problems: max
relative gap \num{1.4e-16}. An independent verifier repeated this with its own
seeds, spacings and neighbour counts over 20 further problems and 60
comparisons, and reported a maximum gap of \textbf{exactly zero}, bit-identical,
while confirming that the guess genuinely changes the number of rows used and
therefore the search path.

\paragraph{The learned predictor did not pay off.} A physics-informed graph
network (\num{892}k parameters) was trained on 320 scenarios across eight
families, split by scenario digest (258/39/23), with a loss combining
imitation of the solver's active set and duals with a physics term evaluating
real constraint residuals through the stored sensitivity tensors. Training is
exactly reproducible: two runs at the same seed gave identical validation loss
(difference \num{0}). Calibrated for recall, it retains \SI{64.6}{\percent} of
rows at \SI{100}{\percent} recall, or \SI{25.5}{\percent} at
\SI{96.2}{\percent}.

On 18 held-out large scenarios the objective gap remained exactly zero --- the
certification is indifferent to the model's quality --- but the median
end-to-end speedup including inference was $0.77\times$, against roughly
$4\times$ for simply starting from the \emph{empty} set. Two measurable
reasons: the active set is far smaller than the network's output (two binding
rows out of 1716, against 2--146 proposed), and inference costs
\SIrange{8}{56}{\milli\second} against a \SIrange{14}{180}{\milli\second}
solve.

A wider search confirmed and sharpened this. Four row-selection strategies were
compared, all returning a bit-identical objective: the empty guess
(\numrange{1.3}{10.6}$\times$, using 0--9 rows of 1716), the trained network
(\numrange{0.16}{1.91}$\times$, 119--199 rows), a \emph{certified} analytic
per-row prune (\numrange{0.65}{0.98}$\times$), and that prune combined with an
analytic a-posteriori guess (\numrange{0.43}{0.97}$\times$). The certified prune
deserves a note: a row can only be violated if the budget can move it past its
nominal margin, which is checkable in closed form, and over 8\,896 rows it lost
\textbf{zero} binding rows while pruning \SIrange{15}{65}{\percent}. It is
sound, and it is still slower than predicting nothing.

The claim that survives is about the mechanism: an active-set guess can be
turned into an exactly-certified speedup that grows with problem size. The claim
that does not survive is that \emph{any} predictor --- learned or analytic ---
supplies a better guess than the trivial one. The active set holds two to eleven
rows out of seventeen hundred, so a useful predictor would have to be cheaper
than free. We report this as measured, and note the general pattern it belongs
to: every acceleration that worked in this system replaced a guess with a cheap
way to \emph{know}.

\paragraph{Triage: a ceiling that means nothing.} A separate bake-off asked
whether a model could decide, from mean elements alone and before any
propagation, that a scenario needs no maneuver at all --- 53\% of a benchmark
sweep needs none and pays full screening cost anyway. Every model family
reached ROC-AUC \num{1.0} and specificity \num{1.0} at recall \num{1.0} on a
held-out split: a two-feature closed-form rule at \SI{0.57}{\micro\second},
logistic regression, gradient boosting, random forests, $k$-NN, MLPs, deep
ensembles, MC-dropout, Bayes-by-backprop, a Bayesian LSTM, DeepSets, a Set
Transformer and a 1D CNN.

That ceiling is an artifact. The label is constant within nine of the ten
scenario families, so the task is family recognition and two pooled statistics
accomplish it. A leave-one-family-out refit makes the failure explicit: with
\texttt{induced-cascade} held out, the rule skips \textbf{100 of 100} scenarios
that needed a maneuver --- recall \num{0.000}. Since the one intolerable error
for this system is a false negative, no triage gate is deployed. The genuinely
sound physics rule --- skip if no pair passes the perigee/apogee prefilter ---
fires on \textbf{0 of 423} scenarios, because every synthetic family lives on a
single \SI{550}{\kilo\meter} shell. A benchmark that could discriminate here
needs families whose label varies \emph{within} the family.

